% 第 4 章：换流器短路贡献
\section{换流器短路贡献}\label{sec:sc-converter}

\subsection{两类换流器短路模型}
换流器短路行为由 \fld{SCConverterModel}（\srcpath{short\_circuit.hpp:88}）区分：
\begin{itemize}
  \item \textbf{非 AC 构网 VSC}：受限电流源。源码先按
        \fld{i\_max\_pu}（非默认 1 优先）、\fld{i\_ac\_max\_pu}、
        \fld{i\_max\_pu}、1 的顺序解析 $k$，再计算
        \begin{equation}
        I_{source}=k\frac{\fld{p\_rated\_mw}}{\sqrt3U_{bus}}.
        \end{equation}
        当 \fld{p\_rated\_mw}$\le10^{-6}$ 时视为缺额定（\texttt{...\_missing\_rating}）、
        贡献保持零（\srcpath{short\_circuit.cpp:1477--1484}）。源到故障点的传递还乘
        $|Z_{k,xfer}|/|Z_k|$ 及两端电流基准换算；
  \item \textbf{AC 构网 VSC}：代码使用电压源阻抗
        \begin{equation}
        z_{conv}=\left(r_{sc}+j\begin{cases}x_{sc},&x_{sc}>10^{-12}\\0.15,&\text{否则}
        \end{cases}\right)
        \frac{S_b}{\begin{cases}p_{rated},&p_{rated}>10^{-6}\\S_b,&\text{否则}.
        \end{cases}}
        \end{equation}
        贡献按 $c/|z_{conv}|$ 和传递阻抗比计算。
\end{itemize}
上述实现见 \srcpath{src/short\_circuit/short\_circuit.cpp:1273--1279}、
\srcpath{short\_circuit.cpp:1321--1342} 和 \srcpath{short\_circuit.cpp:1449--1496}。

\subsection{贡献结果}
逐换流器贡献记于 \fld{SCConverterContributionResult}（:144）：\fld{converter\_index}、\fld{bus\_id}、
\fld{model}（\texttt{grid\_following\_current\_source} / \texttt{ac\_grid\_forming\_voltage\_source}）、
\fld{ac\_grid\_forming}/\fld{dc\_grid\_forming} 标志、\fld{p\_rated\_mw}、\fld{i\_limit\_pu} 与
\fld{contribution\_ka}。该贡献并入故障母线 \fld{ikss\_converter\_contrib\_ka}（:114），与发电机/
电机/外网贡献分列，便于按来源审计。

\subsection{与旋转机的区别}
与同步机（内阻抗 $z_G''$ + 校正 $K_G$，见第~\ref{sec:sc-detailed}~章）不同，电力电子换流器的
短路电流受控制限流主导而非物理次暂态电抗，故跟网型贡献接近额定电流量级而非数倍。构网型
在暂态窗口内表现为受阻抗限制的电压源，其持续贡献取决于内电势与等效阻抗。二者均是稳态/
准稳态故障水平估计，\textbf{不}含控制环动态、FRT 轨迹与限流器切换时序。
